
When @@M0@@ is proper over @@M1@@, it is known [4] that @@M2@@ is representable by a group scheme @@M3@@ locally of finite type over @@M4@@. The neutral component @@M5@@ of @@M6@@ is a connected commutative algebraic group over @@M7@@.

\sgaStableTarget{sga6.e10.def.2-4-1}\textbf{Definition 2.4.1.} \uline{Let} @@M8@@ \uline{be a proper scheme over a field} @@M9@@, \uline{and} @@M10@@ \uline{an invertible} @@M11@@\uline{-module} (resp. @@M12@@ \uline{an element of} @@M13@@\uline{). We say that} @@M14@@ (resp. @@M15@@\uline{) is} @@M16@@\uline{-equivalent to zero when there is an integer} @@M17@@ \uline{such that} @@M18@@ (resp. @@M19@@\uline{) is algebraically equivalent to zero. The elements of} @@M20@@ \uline{algebraically equivalent} (resp. @@M21@@\uline{-equivalent) to zero form a subgroup} @@M22@@ (resp. @@M23@@\uline{) of} @@M24@@.

\sgaStableTarget{sga6.e10.lem.2-4-3}\textbf{Lemma 2.4.3.} \uline{For every proper scheme} @@M25@@ \uline{over an algebraically closed field} @@M26@@, \uline{the group} @@M27@@ \uline{is divisible by every integer} @@M28@@ \uline{prime to the characteristic of} @@M29@@.

\emph{Proof}: Since @@M30@@ is algebraically closed, one has @@M31@@. Since @@M32@@ is connected and of finite type, the @@M33@@-th-power endomorphism is surjective when @@M34@@ is prime to the characteristic of @@M35@@ (SGA 3 XV 1.3 and VI 1.3.2). But @@M36@@ is algebraically closed, so the corresponding endomorphism of the group @@M37@@ is surjective, Q.E.D.

From 2.4.3 one concludes that, in the exact sequence
@@M38@@
@@M39@@ is divisible by a prime number @@M40@@ and @@M41@@ is a torsion group, so one can apply 2.1.3 to the restriction of the polynomial map (2.3.1) to @@M42@@. One thus obtains:

